% ============================================================
% Section 49: 中子非弹性散射与声子色散测量
% ============================================================
\section{固体理论：中子非弹性散射与声子色散测量}
\label{sec:neutron-scattering}

\begin{modelbox}[题目：用中子衍射测量声子色散关系]
中子衍射可用于测量晶体中某种激发的 $\omega$ 和 $\vec k$ 的关系。假设晶体的对称性已知，写出衍射的能量、动量守恒式。指出必须测量什么参数以确定 $\omega(\vec q)$。
\end{modelbox}

\begin{tipbox}[考点快速分析]
\begin{itemize}
    \item \textbf{中子非弹性散射}：中子与声子交换能量和动量，是测量声子色散的``金标准''
    \item \textbf{能量守恒}：中子动能变化 = 声子能量 $\hbar\omega(\vec q)$（吸收或发射）
    \item \textbf{准动量守恒}：模去倒格矢 $\vec G$（Umklapp 过程）
    \item \textbf{实验三要素}：入射波矢 $\vec k$、散射波矢 $\vec k'$、散射方向（即 $\vec Q=\vec k-\vec k'$）
\end{itemize}
\end{tipbox}

\begin{derivationbox}[标准解题过程]
\textbf{Step 1: 物理图像}

中子不带电，不受电子云屏蔽影响，可直接与原子核发生强相互作用。当晶格中存在声子（量子化的晶格振动）时，入射中子可以与声子发生\textbf{非弹性散射}：
\begin{itemize}
    \item \textbf{发射声子}：中子损失能量 $\hbar\omega$，动量转移给一个声子
    \item \textbf{吸收声子}：中子获得能量 $\hbar\omega$，消灭一个已有声子
\end{itemize}

设中子质量为 $M$，入射波矢为 $\vec k$，散射波矢为 $\vec k'$：
\begin{equation}
E = \frac{\hbar^2k^2}{2M}, \qquad E' = \frac{\hbar^2k'^2}{2M}.
\end{equation}

\textbf{Step 2: 能量守恒}

散射前后中子能量变化等于声子能量（含正负号）：
\begin{equation}
\boxed{\frac{\hbar^2k^2}{2M} = \frac{\hbar^2k'^2}{2M} \pm \hbar\omega(\vec q).}
\label{eq:neutron-energy}
\end{equation}
\begin{itemize}
    \item $+$ 号：\textbf{声子发射}（中子损失能量，$E>E'$）
    \item $-$ 号：\textbf{声子吸收}（中子获得能量，$E<E'$）
\end{itemize}

\textbf{Step 3: 准动量守恒}

晶体具有离散平移对称性，因此动量守恒只在模去一个倒格矢 $\vec G$ 的意义下成立：
\begin{equation}
\boxed{\hbar\vec k = \hbar\vec k' \pm \hbar\vec q - \hbar\vec G \quad\text{或}\quad \vec k - \vec k' = \pm\vec q - \vec G.}
\label{eq:neutron-momentum}
\end{equation}

\textit{Umklapp 过程}：即使 $\vec k-\vec k'$ 超出第一布里渊区，加上（或减去）一个倒格矢 $\vec G$ 后，仍可折叠回第一布里渊区找到一个等价的 $\vec q$。这是晶体周期性带来的独特效应。

\textbf{Step 4: 从守恒方程反推声子色散}

将守恒方程改写成``测量量 → 声子参量''的形式：
\begin{align}
    \hbar\omega(\vec q) &= \Delta E = E - E' = \frac{\hbar^2}{2M}(k^2 - k'^2), \label{eq:phonon-energy-from-neutron}\\
    \vec q &= \vec Q + \vec G = (\vec k - \vec k') + \vec G. \label{eq:phonon-q-from-neutron}
\end{align}
（符号约定：$\Delta E>0$ 对应发射声子，$\Delta E<0$ 对应吸收声子。）

\textbf{实验操作逻辑}：
\begin{enumerate}
    \item 选定入射中子能量 $E$ 和散射方向（即固定 $\vec Q=\vec k-\vec k'$）。
    \item 扫描出射能量 $E'$，记录散射强度 $I(E')$。
    \item 当 $E'$ 满足式 \eqref{eq:phonon-energy-from-neutron} 时，中子与声子发生\textbf{共振非弹性散射}，$I(E')$ 出现尖锐峰值。
    \item 峰位给出该 $\vec Q$ 对应的声子能量 $\hbar\omega$；由式 \eqref{eq:phonon-q-from-neutron} 得到声子波矢 $\vec q$。
    \item 改变晶体取向或散射角，测量不同 $\vec Q$ 方向的峰位，逐点描出声子色散曲线 $\omega(\vec q)$。
\end{enumerate}

\textit{为什么出现峰值？} 中子与晶格相互作用的散射截面正比于声子态密度和玻色占据因子。当能量、动量\textbf{同时满足}守恒条件时，相空间允许的最终态数目出现奇点（类似共振），散射强度显著增强。
\end{derivationbox}

\begin{formulabox}[守恒方程汇总]
\begin{center}
\begin{tabular}{cl}
\toprule
守恒律 & 方程 \\
\midrule
能量守恒 & $\displaystyle \frac{\hbar^2k^2}{2M} = \frac{\hbar^2k'^2}{2M} \pm \hbar\omega(\vec q)$ \\[10pt]
准动量守恒 & $\displaystyle \vec k - \vec k' = \pm\vec q - \vec G$ \\
\bottomrule
\end{tabular}
\end{center}
\end{formulabox}

\begin{insightbox}[实验方法与深层理解]
\textbf{1. 为什么选中子而不是 X射线或光子？}

\begin{itemize}
    \item \textbf{能量匹配}：热中子能量 $\sim10\text{--}100\,\mathrm{meV}$，与声子能量同数量级（$k_BT_{\text{room}}\approx25\,\mathrm{meV}$），可发生显著的非弹性散射。
    \item \textbf{X射线/光子}：能量太高（$\sim\mathrm{keV}$），声子能量仅为入射能量的 $10^{-5}$ 量级，能量分辨率不足以分辨。
    \item \textbf{动量匹配}：中子波长 $\sim1\,\text{\AA}$，与晶格常数同数量级，满足 Bragg 条件的动量转移恰好落在布里渊区尺度上。
\end{itemize}

\textbf{2. 实验必须测量的三个参数}

\begin{enumerate}
    \item \textbf{入射中子波矢} $\vec k$（或能量 $E$）
    \begin{itemize}
        \item 通过单色器晶体（如热解石墨、Ge 单晶）选择特定波长
        \item 常用方法：Bragg 衍射 $n\lambda=2d\sin\theta$
    \end{itemize}
    \item \textbf{散射中子波矢} $\vec k'$（或能量 $E'$）
    \begin{itemize}
        \item 通过分析器晶体选择特定散射能量的中子
        \item 能量分辨率决定可分辨的声子线宽
    \end{itemize}
    \item \textbf{散射方向}（动量转移 $\vec Q=\vec k-\vec k'$）
    \begin{itemize}
        \item 转动探测器和样品，改变散射角 $2\theta$
        \item 确定 $\vec Q$ 在倒空间中的位置
    \end{itemize}
\end{enumerate}

\textbf{3. 三轴谱仪（Triple-Axis Spectrometer, TAS）}

TAS 是三台晶体（单色器--样品--分析器）围绕样品转动的装置：
\begin{itemize}
    \item 第一轴（单色器）：选定入射能量 $E$
    \item 第二轴（样品台）：改变样品取向，调节 $\vec Q$
    \item 第三轴（分析器）：选定出射能量 $E'$
\end{itemize}
固定 $E$ 和 $E'$，扫描 $\vec Q$；或固定 $\vec Q$，扫描 $E-E'$（能量扫描），即可逐点测量色散关系。

\textbf{4. Umklapp 过程的实验意义}

如果只测量第一布里渊区内的声子，似乎不需要 Umklapp。但实际上：
\begin{itemize}
    \item 声子色散 $\omega(\vec q)$ 在第一布里渊区内已经完整（周期性重复）
    \item 但实验上中子动量转移 $|\vec Q|$ 可能很大，需要利用 $\vec q=\vec Q+\vec G$ 把大 $\vec Q$ ``折叠''回第一布里渊区
    \item 这使得同一条色散曲线可在不同 $\vec Q$ 位置被重复观测，用于验证和精修
\end{itemize}

\textbf{5. 声子谱的物理信息}

通过中子散射测得的声子色散 $\omega(\vec q)$ 可提取：
\begin{itemize}
    \item \textbf{力常数}：从色散曲线斜率（长波声速）和曲率（近边界行为）
    \item \textbf{晶格稳定性}：若某支色散出现虚频（$\omega^2<0$），晶格不稳定
    \item \textbf{电子-声子耦合}：通过谱线宽度和非谐效应
    \item \textbf{相变信息}：软模（频率趋于零）预示结构相变
\end{itemize}
\end{insightbox}

\begin{thinkbox}[考场速记：中子散射问题的答题模板]
遇到``用中子/光子散射测声子/磁振子色散''类问题：
\begin{enumerate}
    \item \textbf{写出能量守恒}：$E_i = E_f \pm \hbar\omega(\vec q)$（$+$ 发射，$-$ 吸收）
    \item \textbf{写出准动量守恒}：$\vec k_i - \vec k_f = \pm\vec q - \vec G$
    \item \textbf{说明测量量}：入射能量 $E$、散射能量 $E'$、散射角 $2\theta$（即 $\vec Q$）
    \item \textbf{说明 Umklapp}：$\vec G$ 使大 $\vec Q$ 折叠回第一布里渊区
    \item \textbf{说明为什么选中子}：能量/动量与声子匹配（X射线能量太高，光子动量太小）
\end{enumerate}
\textit{核心口诀}：能量守恒看声子能量，准动量守恒看声子波矢；Umklapp 让大 $\vec Q$ 变小 $\vec q$。
\end{thinkbox}
